%%
%% PLDI 2027 submission: "Transients Without Tears".
%% Automatic in-place conversion of persistent data structure updates.
%%
\documentclass[sigplan,review,anonymous]{acmart}

\usepackage{listings}
\usepackage{booktabs} %% For formal tables
\usepackage{seqsplit}

\lstset{
  language=Lisp,
  basicstyle=\small\ttfamily,
  breaklines=true,
  frame=single,
  numbers=left,
  numberstyle=\tiny,
  keywordstyle=\color{blue},
  commentstyle=\color{green!40!black},
  stringstyle=\color{red},
  showstringspaces=false
}

%%
%% \BibTeX command to typeset BibTeX logo in the docs
\AtBeginDocument{%
  \providecommand\BibTeX{{%
    \normalfont B\kern-0.5em{\scshape i\kern-0.25em b}\kern-0.8em\TeX}}}

%% Rights management information.
\setcopyright{acmlicensed}
\copyrightyear{2027}
\acmYear{2027}
\acmDOI{XXXXXXX.XXXXXXX}

%% These commands are for a PROCEEDINGS abstract or paper.
\acmConference[PLDI '27]{48th ACM SIGPLAN Conference on Programming Language Design and Implementation}{June 2027}{TBD}
\acmBooktitle{Proceedings of the 48th ACM SIGPLAN Conference on Programming Language Design and Implementation (PLDI '27)}
\acmPrice{15.00}
\acmISBN{978-1-4503-XXXX-X/27/06}

%%
%% The "title" command has an optional parameter for a short title.
\title[Transients Without Tears]{Transients Without Tears: Automatic In-Place Conversion of Persistent Data Structure Updates in a Dynamic Language}

\author{Anonymous Author(s)}
\affiliation{%
  \institution{Anonymous Institution}
  \city{Anonymous City}
  \country{Anonymous Country}
}
\email{anonymous@example.com}

\renewcommand{\shortauthors}{Anonymous Author(s)}

\begin{abstract}
Persistent-first languages pay a large allocation tax on accumulation
loops. Languages like Koka eliminate it with ownership types (Perceus/FBIP);
Clojure asks the programmer to opt in manually with transients. We show the
optimization can be made *automatic and sound* in a latently-typed,
open-world language with live redefinition---with no type system---via (1) a
tail-position-sensitive *usage classification* that proves the *linear usage*
of loop and fold accumulators, (2) a rewriter aligned
with the classifier *by construction*, targeting edit-tagged in-place
transients, and (3) *world-guarded regions* that keep the optimization sound
when definitions change at runtime. The result: 4.1--7.6x on
accumulation-bound code, byte-identical outputs on a production-scale event
simulator, zero measured cost when disabled or inapplicable, and an honest
cost model of when transient conversion pays.
\end{abstract}

\begin{CCSXML}
<ccs2012>
   <concept>
       <concept_id>10011007.10010940.10010941.10010942.10010943</concept_id>
       <concept_desc>Software and its engineering~Compilers</concept_desc>
       <concept_significance>500</concept_significance>
       </concept>
   <concept>
       <concept_id>10011007.10011006.10011041.10011045</concept_id>
       <concept_desc>Software and its engineering~Dynamic compilers</concept_desc>
       <concept_significance>500</concept_significance>
       </concept>
   <concept>
       <concept_id>10003752.10003753.10003757</concept_id>
       <concept_desc>Theory of computation~Program analysis</concept_desc>
       <concept_significance>500</concept_significance>
       </concept>
 </ccs2012>
\end{CCSXML}

\ccsdesc[500]{Software and its engineering~Compilers}
\ccsdesc[500]{Software and its engineering~Dynamic compilers}
\ccsdesc[500]{Theory of computation~Program analysis}

\keywords{escape analysis, persistent data structures, transients, in-place update, dynamic languages, lisp, speculative optimization, deoptimization}

\begin{document}
\maketitle

\section{Introduction}

Persistent data structures impose a significant performance tax. In FOL, our
Lisp dialect, accumulation-heavy benchmarks can run up to 20$\times$ slower and
allocate **900$\times$** more memory than their mutable counterparts, dominated by
garbage collection pressure. The canonical pattern is an accumulator in a loop,
such as \texttt{(loop [acc \{\}] ... (recur (assoc acc k v)))} or a
\texttt{reduce}, where each update allocates a new version of the collection.

Existing solutions do not transfer to our setting. Perceus/FBIP
\cite{reinking2021perceus} requires ownership types, which FOL lacks. Clojure's
transients are a manual, unsafe escape hatch. Crucially, FOL is a dynamic,
open-world language where method and class definitions can change while
optimized code is live, precluding any static analysis that assumes a closed
world.

This paper shows how to make transient conversion automatic and sound in this
hostile environment. Our contributions are:

\begin{enumerate}
    \item \textbf{Linear usage analysis.} A flow- and tail-position-sensitive
    analysis that proves an accumulator is used linearly---that is, uniquely at
    the point of update---and can therefore be safely converted to an in-place
    transient, without requiring a formal type system.
    \item \textbf{Alignment-by-construction.} The analysis and transformation are
    a single pass, ensuring that only provably safe patterns are rewritten.
    \item \textbf{Name-resolution-faithful summaries.} A table of standard
    library operator effects that mirrors the compiler's own name resolution,
    ensuring the analysis is as sound as the language's semantics.
    \item \textbf{Soundness under live redefinition.} A world-guard mechanism
    emits dual-path code that safely falls back to the persistent path if a
    function or class assumed by the optimization is redefined at runtime.
    \item \textbf{An honest empirical cost model.} We evaluate the optimization on
    microbenchmarks and a real-world event simulator, quantifying both the wins
    (4-7x) and the limitations, including a motivating benchmark that our
    intra-procedural analysis cannot convert.
\end{enumerate}

\section{Background and Motivating Example}

FOL is a Lisp-1 dialect with persistent vectors (32-way tries), HAMT
dictionaries, and a transpiler to Common Lisp. Standard library functions are
resolved by name at emit time.

Consider the \texttt{lh-pop-batch} function from our event simulator benchmark.
It drains a queue, building up a result vector with \texttt{conj}.

\begin{lstlisting}
;; --- Source ---
(defn lh-pop-batch [q n]
  (loop [acc [] q q i 0]
    (if (or (>= i n) (empty? q))
      [acc q]  ; acc exits inside a literal
      (recur (conj acc (peek q)) (pop q) (inc i)))))

;; --- After transient conversion ---
(let [validity-cell (register-region '("conj" "peek" "pop"))]
  (if (car validity-cell)
    ;; Fast path
    (loop [acc (transient []) q q i 0]
      (if (or (>= i n) (empty? q))
        [(persistent! acc) q]
        (recur (conj! acc (peek q)) (pop q) (inc i))))
    ;; Fallback path (original code)
    ...))
\end{lstlisting}

Our system automatically converts this loop to use a transient accumulator. The
\texttt{conj} becomes a mutating \texttt{conj!}, and the accumulator is sealed
with \texttt{persistent!} at the exit. The entire fast path is wrapped in a
world guard that checks if any of the assumed functions have been redefined.

\section{The Analysis}

Our analysis proves **linear usage**, a property analogous to that enforced by
linear type systems \cite{wadler1990linear} but established here by a syntactic
usage analysis. A value is used linearly if it is fresh, unaliased, and every
use is a last-use that flows into a valid update or exit.

\subsection{Tier-1 Summaries}

We use a hand-verified table of ~75 standard library function summaries. Each
summary describes the function's effect on its arguments using a simple lattice:
\texttt{:none < :invoked < :shared-with-result < :retained}. A key assumption is
the **root-level convention**: persistent operations like \texttt{assoc} produce
a fresh root node, which is safe for all clients because our transient protocol
is copy-on-write.

\subsection{Name-Resolution Faithfulness}

The analysis is only as sound as its knowledge of what function is being
called. Our summary table mirrors the compiler's own name resolution logic,
including per-file lexical definitions. This leads to a key lemma: any program
the analysis misjudges is a program the compiler itself would have resolved the
same way.

\subsection{The Classifier}

The core of the analysis is a recursive walk over the AST that classifies every
use of a candidate accumulator. It propagates tail-position context through
\texttt{if}, \texttt{do}, \texttt{bind}, etc. It recognizes sanctioned uses
like being the argument to an update function (\texttt{conj}, \texttt{assoc})
or appearing at a valid exit point. Any other use, such as being stored in
another data structure or captured in a closure, disqualifies the accumulator.

\subsection{Chains}

The analysis recognizes chains of linear updates, including direct calls,
\texttt{->} threading macros, and if-branched chains. It also handles **reduce
init-threading**: a \texttt{reduce} whose lambda is proven to be linear becomes
a valid link in an update chain itself.

\subsection{Formal Core}

We formalize the analysis for a core calculus. The key judgment, `Γ; Σ ⊢ e ⇓ U`,
states that expression `e` has usage `U` in context `Γ` given operator summaries
`Σ`. We show that "messy" implementation features are sound elaborations of this
core: `->` threads and `if`-branched chains desugar into nested applications
that the core rules handle, and name-resolution faithfulness is a property of
how `Σ` is constructed, not the rules themselves. The main theorem (proof
sketched in an appendix) states that if an accumulator qualifies, the
transient-converted program is observationally equivalent to the original,
proven via bisimulation against a store semantics. The termination and
correctness of the inter-procedural fixpoint analysis are discussed formally in
Appendix A.

\section{Transformation and Representation}

The transformation is aligned by construction with the analysis.

\subsection{Rewriter}

If an accumulator qualifies, the rewriter wraps its initializer with
\texttt{transient}, replaces all update calls with their `!` counterparts (e.g.,
\texttt{conj} -> \texttt{conj!}), and inserts \texttt{persistent!} at all valid
exit points.

\subsection{Edit-Tagged Transients}

Our implementation uses edit-tagged transients, a token-based copy-on-write
(COW) HAMT for dictionaries and a vector with an owned tail. This is crucial as
it allows reads on transients, unlike simpler wrapper-based approaches. A plain
persistent lookup would silently misread edited data; our transient-aware read
primitives correctly handle lookups on the mutable view.

\section{Soundness under Live Redefinition}

The primary threat to soundness is live redefinition. If our optimization bakes
in the meaning of \texttt{assoc}, a later \texttt{(defn assoc ...)} could
invalidate the compiled code.

Our solution is a dual-path emission strategy. At load time, each optimized
region registers its dependencies (e.g., `{"assoc", "conj"}`) and receives a
validity cell. At execution time, defining a function like \texttt{assoc} calls
\texttt{note-redefinition}, which looks up all dependent regions and flips their
validity cells. The compiled code checks this cell on entry.

\begin{lstlisting}
(IF (CAR (LOAD-TIME-VALUE (REGISTER-REGION '("assoc"))))
    ;; Fast path: use assoc!
    ...
    ;; Fallback path: use original persistent assoc
    ...)
\end{lstlisting}

This mechanism is a static, region-scoped analogue to a JIT's world-age or
deoptimization system, but requires no runtime recompilation.

\section{Implementation}

The system is implemented in ~2,300 lines of Common Lisp across the analysis,
rewriter, and world-guard machinery. It includes an "audit mode" that runs the
analysis without emitting code. With only intra-procedural (Tier-1) analysis,
the audit found 9 convertible candidates out of 17 loops in our LSim benchmark.
The DVI benchmark, a key motivator whose accumulator is passed through a user
function, had 0 candidates. With the addition of our inter-procedural (Tier-2)
summary inference engine, the DVI benchmark's helper function is correctly
summarized, and its main loop now converts successfully.

\section{Evaluation}

We evaluate correctness, performance on microbenchmarks, and end-to-end impact.
All benchmarks were run on a Ryzen 9 5900X with SBCL 2.6.0.

\textbf{RQ1 (Correctness):} Our full test suite of 3,248 checks passes with the
optimization enabled. The LSim event simulator produces byte-identical output.

\textbf{RQ2 (Speedup):} On accumulation-bound microbenchmarks, we see speedups
of 4.1x to 7.6x and memory allocation reductions of 8-11x.

\begin{table}[h]
  \caption{Microbenchmark results (best-of-3).}
  \label{tab:speedup}
  \small
  \begin{tabular}{lrrrr}
    \toprule
    Workload & Baseline & Optimized & Speedup & Mem. Reduction \\
    \midrule
    dict loop (200k) & 224.6ms & 55.2ms & **4.07×** & 8.7x \\
    vector loop (1M) & 309.9ms & 41.0ms & **7.56×** & 11.0x \\
    reduce (500k) & 324.1ms & 48.8ms & **6.64×** & 9.1x \\
    DVI (200k) & 241.2ms & 63.1ms & **3.82×** & 8.5x \\
    \bottomrule
  \end{tabular}
\end{table}

\textbf{RQ3 (End-to-end):} On the LSim simulator, the wall-clock speedup is
modest (1.04x), though allocation is reduced by 1.10x. This is explained by our
cost model: the hot loops in LSim perform only small-N accumulations, where the
constant-factor overhead of the transient boundaries limits the win. The
dominant cost is heap churn from other sources, which our technique does not
target.

\textbf{RQ4 (Cost):} When disabled or inapplicable, the performance is identical
to the baseline within measurement noise. The world guards are effectively free.

\section{Discussion and Limitations}

Our inter-procedural analysis significantly broadens the optimization's reach,
as demonstrated by the successful conversion of the DVI benchmark. However,
limitations remain.

\textbf{Inter-procedural Analysis Scope.} Our summary inference engine operates
on a file-at-a-time basis. While it can analyze through any chain of calls
within a single compilation unit, it cannot see across module boundaries. A
function imported from another module is treated as a black box, and passing an
accumulator to it will disqualify the conversion. A whole-program analysis or a
system of explicit, checked interface contracts would be required to lift this
limitation.

\textbf{Profitability Heuristics.} As our end-to-end evaluation on LSim shows,
the performance gain on small-N accumulations is bounded by the constant-factor
overhead of creating and sealing transient collections. For loops with very few
iterations, this can lead to a performance regression. A profitability
heuristic could prevent this by augmenting the world guard with a runtime check,
such as `(> (count accumulator) <threshold>)`, where the threshold is an
empirically determined constant. This would ensure the optimization is only
applied when likely to be beneficial, making it safer to enable by default.

Finally, our transient representation for sets does not yet support reads. This
is an implementation limitation: due to the unordered nature of the underlying
hash table, providing a consistent copy-on-write view for reads is more complex
than for indexed collections. This restricts the optimization's applicability
for set-based accumulators in read-heavy loops.

\section{Related Work}

Our work is an alternative to type-system-based approaches like Perceus/FBIP
\cite{reinking2021perceus,lorenzen2023functional} for languages that lack
ownership or linear types, and automates the manual process of using Clojure's
transients \cite{hickey2009clojure}. Unlike classic escape analysis
\cite{choi1999escape,blanchet2003escape}, which is a *may-analysis* for
stack allocation, our linear usage analysis is a *must-analysis* for proving
the safety of in-place mutation.

Our world-guard mechanism is a static, ahead-of-time (AOT) analogue to the
assumption-and-deoptimization systems used in Just-in-Time (JIT) compilers. The
choice of a static approach is dictated by FOL's nature as a transpiled AOT
language, which lacks the runtime infrastructure for JIT-style dynamic
recompilation. A JIT might speculatively optimize a loop assuming a function is
not redefined, and then deoptimize back to an interpreter if that assumption is
violated. Our approach, in contrast, emits both the fast and fallback paths
ahead of time, selecting between them with a simple guard.

This AOT strategy has distinct advantages: it avoids the complexity and runtime
overhead of a JIT, including on-stack replacement machinery and interpreter
frame metadata. The performance is predictable, with no stalls from runtime
recompilation. The primary trade-off is granularity: our invalidation occurs at
region entry, whereas a JIT can deoptimize a running activation mid-loop. Our
system is most similar to Julia's world-age mechanism \cite{belyakova2020oopsla},
but we apply it to uniqueness and escape facts rather than method dispatch tables.

\section{Conclusion}

We have presented a system for automatically converting allocation-heavy loops
over persistent data structures to use in-place transient updates. Our approach is
sound in a dynamic, open-world language, combining a syntactic linear usage
analysis with a world-guard mechanism to handle live redefinition. The result
is a significant performance improvement for accumulation-bound code, achieved
without programmer annotation or a formal type system.

\bibliographystyle{ACM-Reference-Format}
\bibliography{pldi2027}

\appendix
\section{Fixpoint Analysis Formalism}

Our inter-procedural analysis for inferring function summaries (Tier-2)
operates by iterating to a fixpoint. This appendix provides a formal sketch of
this process, establishing its termination and correctness.

\subsection{The Summary Lattice}

The analysis computes a usage summary for each function. A summary for a
function with $k$ parameters is a tuple $(u_1, u_2, \dots, u_k)$, where each
$u_i$ is an element from the finite usage lattice $\mathcal{U}$:

\[ \mathcal{U} = \{\texttt{:none} \sqsubset \texttt{:invoked} \sqsubset \texttt{:shared-with-result} \sqsubset \texttt{:retained}\} \]

Here, $\sqsubset$ denotes "is more precise than". The top element,
\texttt{:retained}, is the most conservative assumption (the parameter may be
retained indefinitely), while \texttt{:none} is the most precise (the parameter
is not used). The lattice of summaries for a single function is the product
lattice $\mathcal{U}^k$, which is also finite.

\subsection{Transfer Function and Fixpoint Iteration}

Let $S$ be the set of all function summaries in a compilation unit (a file). We
define a transfer function $F: S \to S$ that computes a new set of summaries
based on the current ones. For each function $f$ in the unit, $F$ re-analyzes
its body using the current summaries for all functions it calls.

The analysis proceeds by iterating $S_{i+1} = F(S_i)$, starting with the most
conservative assumption, $S_0$, where all parameters of all user-defined
functions are marked as \texttt{:retained}.

\textbf{Termination.} The analysis is guaranteed to terminate because:
\begin{enumerate}
    \item The overall summary space for the compilation unit is finite, as it is
    the product of the finite summary lattices for each function.
    \item The transfer function $F$ is monotonic. As the analysis gains more
    precise information about a callee (i.e., its summary moves down the
    lattice), the inferred summary for the caller can only become more precise
    or stay the same. It can never become less precise (move up the lattice).
\end{enumerate}

Since the iteration proceeds over a finite lattice with a monotonic function, it
is guaranteed to reach a fixpoint in a finite number of steps. In practice, for
typical call graphs, this fixpoint is reached very quickly (2-3 iterations).

\end{document}